% Grove of Knowledge v6 --- clean publication skeleton
% Pilot results preserved as comments (% PILOT: ...)
% [TBD] marks cells awaiting publication experiments

\documentclass{article}
\usepackage[preprint,nonatbib]{neurips_2024}
\usepackage[utf8]{inputenc}
\usepackage[T1]{fontenc}
\usepackage{lmodern}
\usepackage{hyperref}
\usepackage{url}
\usepackage{booktabs}
\usepackage{amsfonts}
\usepackage{amsmath}
\usepackage{nicefrac}
\usepackage{microtype}
\usepackage{graphicx}
\usepackage{xcolor}
\usepackage{enumitem}

\newcommand{\grove}{Grove of Knowledge}

\title{Grove of Knowledge: Self-Organizing Expert Trees for the \\ Long Tail of Language Model Knowledge}

\author{
  Erik de Bruijn \\
  Independent Researcher \\
  \texttt{erik@erikdebruijn.nl}
}

\begin{document}
\maketitle

% ============================================================================
\begin{abstract}
Language models encode all knowledge---from common syntax to rare domain
terminology---in shared parameters. Adding knowledge means adding
parameters, and therefore inference cost. We propose the \grove{}: a
frozen base model augmented with hot-pluggable, per-layer gated LoRA
adapters that grow where needed and cost nothing when inactive.

The frozen base serves as a permanent epistemic reference: because its
weights never change, we can always measure what an adapter adds by
comparing expert loss to base loss on any token. This \emph{relative
learntropy}---closely related to Schmidhuber's compression
progress---becomes the unified signal for \emph{what} to learn (token
selection), \emph{when} to stop (Goldilocks zone detection), and
\emph{where} to grow (adapter spawning).

A per-layer delta gate controls each adapter's contribution:
$y = f_{\text{base}}(x) + \sigma(g_l) \cdot (f_{\text{adapter}}(x) -
f_{\text{base}}(x))$, initialized near OFF ($\sigma(-2) \approx 0.12$).
The gate learns from data which layers need adaptation and which should
pass through unchanged, protecting generic capability while injecting
domain knowledge. On Qwen3-8B, standard PEFT LoRA degrades generic
perplexity by 25--32\%; our gated approach \emph{improves} it by 6\%.
Ten independently-trained adapters compose via joint gate fine-tuning
with 94\% of composition checks passing.

The grove builds its own curriculum: scoring candidate data by expected
compression progress, routing to the adapter with highest learning
potential, and spawning new adapters when existing ones plateau. The
result is a self-organizing knowledge system where the long tail of
domain knowledge lives in lightweight, composable specialists---each
measured in kilobytes---while the shared trunk preserves general
reasoning ability intact.
\end{abstract}

% ============================================================================
\section{Introduction}
\label{sec:intro}

Knowledge in language models follows a power-law distribution. A small
core of common patterns (syntax, high-frequency facts, conversational
conventions) accounts for the majority of token predictions, while a
long tail of specialized knowledge (rare medical conditions, niche
programming APIs, regional legal codes) accounts for the majority of
\emph{distinct} knowledge. Current models embed both in shared
parameters. The only way to know more is to be larger.

This coupling between knowledge breadth and model size creates a
practical problem: every domain a model covers increases its inference
cost for \emph{all} queries, including those that need none of that
domain knowledge. Retrieval-augmented generation sidesteps parameter
growth but adds context tokens. Tool use adds generation tokens. Both
increase per-query latency proportional to the knowledge accessed.

We propose an alternative: place the long tail in separately-loadable
adapters whose inference cost is \emph{zero} until activated. The dense
base model handles general reasoning and is always resident in memory.
Domain knowledge lives in lightweight LoRA adapters (4\,KB--256\,KB
each) that merge into the forward pass on demand. Adding a thousand
specialists does not slow a single generic query.

The key insight enabling this architecture is that a \emph{frozen} base
model is not merely a computational convenience---it is an epistemic
reference point. Because the base weights never change, we can always
compute what any adapter adds: the difference between expert loss and
base loss on any token is a direct measure of the adapter's knowledge
contribution. This relative signal---which we connect to Schmidhuber's
compression progress~\cite{schmidhuber2008} and Wozniak's
learntropy~\cite{wozniak2018learntropy}---drives all adaptation
decisions in the grove.

A per-layer delta gate protects the base model's generic capability.
Rather than applying adapter modifications uniformly (as in standard PEFT
LoRA, which degrades generic perplexity by 25--32\%), the gate learns
from data which layers need domain-specific overrides and which should
pass through unchanged. The result: domain knowledge injection with
simultaneous \emph{improvement} of generic perplexity.

\paragraph{Contributions.}
\begin{enumerate}[nosep,leftmargin=1.5em]
\item \textbf{Relative learntropy as a learning signal.} We formalize
      the frozen base as an epistemic reference, enabling a
      compression-progress signal that drives token selection, training
      duration, and adapter spawning (Section~\ref{sec:learning}).
\item \textbf{Per-layer delta gates.} A learned gate per layer controls
      adapter contribution, protecting generic capability while enabling
      selective knowledge injection. Standard PEFT degrades generic PPL
      by 25--32\%; gated adapters improve it by 6\%
      (Section~\ref{sec:architecture}).
\item \textbf{Self-organizing grove growth.} Adapters spawn, split, and
      prune based on compression progress signals, building their own
      curriculum from a multi-domain data pool
      (Section~\ref{sec:learning}).
\item \textbf{Composable multi-adapter deployment.} Ten
      independently-trained adapters compose via joint gate fine-tuning,
      with 94\% of composition checks passing and 77\% throughput gain
      from sparse routing (Section~\ref{sec:architecture}).
\item \textbf{Knowledge PPL as an evaluation metric.} Splitting
      evaluation tokens by base-model surprisal reveals that adapters
      improve knowledge-token perplexity by 64.9\%, a gain hidden by
      uniform PPL (Section~\ref{sec:architecture}).
\end{enumerate}

% ============================================================================
\section{Learning by Active Compression}
\label{sec:learning}

Before describing the architecture, we describe how the grove learns.
The central claim is that the frozen base model transforms adapter
training from passive parameter fitting into active, directed
compression---and that this active learning is more naturally described
as compression progress than as loss minimization.

% ----------------------------------------------------------------------------
\subsection{Compression Progress as the Learning Signal}
\label{sec:compression-progress}

Schmidhuber~\cite{schmidhuber2008} argues that curiosity---the
intrinsic motivation to learn---is driven not by prediction error itself,
but by \emph{improvement in prediction}: the first derivative of the
learner's compression ability. A TV showing static has high prediction
error but zero compression progress; it is maximally surprising yet
completely uninteresting. A well-paced lecture has moderate prediction
error but high compression progress; it is the productive zone.

This distinction is critical for adapter training. Raw cross-entropy
loss measures how badly the model predicts a token, but it does not
distinguish between tokens that are \emph{productively hard} (the model
can learn from them) and tokens that are \emph{irreducibly hard} (noise,
out-of-distribution artifacts, or patterns beyond the adapter's
capacity). Training on all high-loss tokens equally wastes gradient
budget on the irreducible portion.

Wozniak's learntropy~\cite{wozniak2018learntropy} captures this
distinction from a cognitive science perspective: learntropy is the
\emph{attractiveness} of a signal for a specific learner at a specific
moment. It depends on the learner's prior knowledge, can go negative
(when the learning attempt causes interference), and follows an
inverted-U relationship with signal difficulty---the Goldilocks
zone~\cite{vygotsky1978}.

We operationalize this insight in the adapter training loop. Learntropy
in our framework is \emph{not} raw cross-entropy. It is the rate at
which an adapter's compression ability improves on a given data stream,
measured relative to the frozen base.

% ----------------------------------------------------------------------------
\subsection{Relative Learntropy: Unique to Frozen-Base Architectures}
\label{sec:relative-learntropy}

Standard fine-tuning overwrites the base model's weights, destroying the
reference point needed to measure what has been learned. In our
architecture, the base model is frozen: its loss on any token is always
available. This enables a \emph{relative} learning signal.

For a token $t$ in context $c$, the base model produces loss
$\ell_{\text{base}}(t) = -\log P_{\text{base}}(t \mid c)$ and the
adapter-augmented model produces
$\ell_{\text{expert}}(t) = -\log P_{\text{expert}}(t \mid c)$.
The relative learntropy of token $t$ for expert $e$ is:
\begin{equation}
\label{eq:relative-learntropy}
    \Delta\ell(t, e) = \ell_{\text{expert}}(t) - \ell_{\text{base}}(t)
\end{equation}
This measures what the adapter adds (or subtracts) relative to the
base. Negative $\Delta\ell$ means the adapter has learned something the
base does not know. Positive $\Delta\ell$ means the adapter is making
things worse---a signal to stop training on this token type.

The connection to Rho-1~\cite{rho1_2024} is direct. Rho-1 scores each
pretraining token by the gap between a training model's loss and a
reference model's loss, training only on tokens where the gap exceeds a
threshold. They achieve 30\% absolute improvement on math benchmarks
using only 3\% of pretraining tokens. Our architecture provides the same
signal for free: the frozen base \emph{is} the reference model, and the
gap is computable at every training step without additional forward
passes (the base forward pass is already needed for the gated
architecture).

% PILOT: Relative surprise (L_expert - L_base) correlation with raw CE: 0.64 (not equivalent)
% PILOT: Detects overfitting onset at step 1600 (generic relative surprise drops sharply)

The first derivative of relative learntropy over training steps---how
fast the adapter is improving on this data---is Schmidhuber's
compression progress applied to the adapter:
\begin{equation}
\label{eq:compression-progress}
    \text{CP}(e, \mathcal{D}, s) =
    -\frac{d}{ds} \mathbb{E}_{t \in \mathcal{D}}
    \bigl[\Delta\ell(t, e)\bigr]
\end{equation}
Positive compression progress means the adapter is actively learning
from data $\mathcal{D}$ at step $s$. Zero compression progress means it
has plateaued---either it has extracted all learnable information (mastery)
or the data is beyond its capacity (noise). Negative compression
progress means it is forgetting.

% PILOT: Learning speed signal finds Goldilocks zone at steps 300-2950
% PILOT: Detects overfitting 50 steps before relative surprise

% ----------------------------------------------------------------------------
\subsection{The Active Learning Cycle}
\label{sec:active-cycle}

A candidate data stream must pass through a two-stage funnel before
the grove invests GPU time in training:

\paragraph{Stage 1: Compressibility filter.} Is the data structured at
all? A fast CPU-side check (e.g., gzip ratio, BPE tokenizer efficiency)
rejects data that is genuinely random---encrypted blobs,
\texttt{/dev/urandom}, corrupted files. This is a cost optimization, not
an epistemic judgment: it prevents wasting GPU minutes on data that a
gzip call could have filtered in microseconds. Structured data that the
grove cannot currently learn from (e.g., Chinese text when no Chinese
adapter exists) passes this stage---it has compressible structure, even
if no current adapter can exploit it.

\paragraph{Stage 2: Learntropy probe.} The real work. Train each
candidate adapter for a short probe window ($\sim$100 steps) on the
data and measure compression progress. This is inherently
learner-relative: the same dataset may produce high learntropy for
adapter $e_1$ (novel, learnable patterns) and zero learntropy for
adapter $e_2$ (already mastered or beyond capacity). Data that no
adapter can learn from is \emph{parked}---it is not discarded (it may
become learnable when a new adapter is spawned or an existing one
grows), but it does not consume training budget now.

The grove then uses compression progress to build its own curriculum:

\begin{enumerate}[nosep,leftmargin=1.5em]
\item \textbf{Score.} For each data stream $\mathcal{D}_k$ that passed
      the funnel, estimate compression progress
      $\text{CP}(e_i, \mathcal{D}_k)$ for each adapter.
\item \textbf{Route.} Assign $\mathcal{D}_k$ to the adapter with
      highest compression progress. If no adapter has
      positive CP, flag for new adapter creation.
\item \textbf{Train.} Train the selected adapter for $N$ steps.
\item \textbf{Evaluate.} Did compression progress occur?
      If yes, continue. If CP is zero or negative
      for $M$ consecutive evaluations, stop and switch streams.
\item \textbf{Grow.} If a flagged stream shows reducible loss
      (high $\Delta\ell$ that decreases over a probe), spawn a new
      adapter.
\end{enumerate}

% PILOT: Active selection beats random by 6.4pp on primary domain
% PILOT: Goldilocks zone found at steps 300-2950
% PILOT: Knowledge PPL -64.9% on knowledge tokens, +11.5% on style tokens

This cycle is not yet tested at scale (Section~\ref{sec:limitations}),
but pilot experiments demonstrate its components individually: active
token selection outperforms random by 6.4 percentage points on primary
domain perplexity, and the Goldilocks zone (the window of productive
training steps) is detectable from the compression progress signal.

% ============================================================================
\section{Architecture}
\label{sec:architecture}

This section describes the machinery that enables Section~\ref{sec:learning}.
The architecture has three components: per-layer delta gates, a tree
structure for organizing adapters, and growth mechanics for expanding the
grove.

% ----------------------------------------------------------------------------
\subsection{Per-Layer Delta Gate}
\label{sec:delta-gate}

Each adapter layer is controlled by a scalar gate $g_l$ that modulates
the adapter's contribution:
\begin{equation}
\label{eq:delta-gate}
    y = f_{\text{base}}(x) + \sigma(g_l(x)) \cdot
    \bigl(f_{\text{adapter}}(x) - f_{\text{base}}(x)\bigr)
\end{equation}
where $g_l$ is a learned linear projection with bias initialized to
$-2.0$ ($\sigma(-2) \approx 0.12$, starting near OFF). The gate
operates on the \emph{delta}---the difference between adapter and base
output---not on the adapter output directly.

Training is two-phase:
\begin{enumerate}[nosep,leftmargin=1.5em]
\item \textbf{Adapter phase} ($\sim$500 steps): Train LoRA weights on
      domain data with gates frozen at initialization.
\item \textbf{Gate phase} ($\sim$1500 steps): Freeze LoRA weights, train
      gates on mixed domain + generic data. Gates learn which layers
      need adaptation for this domain.
\end{enumerate}

% PILOT: Learntropy-guided recipe: 500+1500 wins both dimensions
% PILOT: Domain -32.1%, generic -12.1%
% PILOT: Gate adds 3.4pp generic protection on top of early stopping

The gate provides a Piagetian interpretation~\cite{piaget1952}: low gate
values indicate \emph{assimilation} (the base model's representation is
sufficient; the adapter acts as a lens redirecting existing computation)
while high gate values indicate \emph{accommodation} (the adapter
overrides the base with structurally new computation). Per-layer gate
magnitudes correlate with domain PPL impact (Spearman $\rho = 0.717$,
$p = 8 \times 10^{-5}$).

The protection effect is robust: standard PEFT LoRA (without gating)
degrades generic perplexity by 25--32\%, while our gated approach
improves it by 6\%. Even a naive output-level gate reduces degradation
by 11$\times$ (from +25\% to +2.9\%), but per-layer gating is strictly
superior.

% PILOT: PEFT LoRA degrades generic PPL +25-32%, gated -6-13%
% PILOT: Gate bias init robust: 4 values (-1 to -4), selectivity 0.607-0.632
% PILOT: 8 gate combinations tested, per-layer delta gate best on both dimensions

% ----------------------------------------------------------------------------
\subsection{Tree Structure}
\label{sec:tree-structure}

Adapters are organized as a tree rooted at the base model. Each adapter
applies LoRA modifications to FFN layers (\texttt{gate\_proj} and
\texttt{up\_proj}), targeting the layers where domain-specific knowledge
lives~\cite{geva2021transformer}. Attention layers remain shared,
preserving general reasoning.

On Qwen3-8B (36 layers), we partition adapter layers into two functional
zones:
\begin{itemize}[nosep,leftmargin=1.5em]
\item \textbf{Identity layers (L0--11):} ``Who am I''---style, register,
      formatting conventions. Low-rank adapters suffice (lens-like).
\item \textbf{Expert layers (L12--35):} ``What do I know''---domain
      facts, specialized vocabulary, technical patterns. Higher-rank
      adapters encode genuinely new knowledge (crystal-like).
\end{itemize}

Sub-experts are created by splitting: when an adapter's compression
progress stalls on a data stream it was previously learning from, the
adapter is a candidate for splitting into two children that can
specialize further.

% ----------------------------------------------------------------------------
\subsection{Growth Mechanics}
\label{sec:growth}

\paragraph{Speculative splitting with rollback.}
When the compression progress signal suggests a split, we create a
child adapter and train it for a probe window. If the child achieves
positive compression progress that the parent could not, the split is
accepted. Otherwise, it is rolled back. In pilot experiments, 2 of 8
attempted splits were accepted and 6 were rolled back---the system
is conservative by design.

% PILOT: Speculative split: 2/8 kept, 6 rolled back (one run)
% PILOT: 8-expert tree: all 8 active, M_ij CV up to 2.34

\paragraph{Joint gate training for multi-adapter composition.}
When multiple adapters are deployed simultaneously, their gates are
fine-tuned jointly on mixed data. This resolves conflicts (two adapters
claiming the same layer) and enables cooperative specialization. Ten
independently-trained adapters compose via this mechanism with 94\% of
composition checks passing; the 6\% of failures are semantically
justified overlaps (physics$\leftrightarrow$chemistry,
biology$\leftrightarrow$medicine).

\paragraph{Sparse routing.}
At inference time, the router identifies which adapters are relevant
(gate activations above threshold) and skips all others. This yields
77\% throughput improvement over dense evaluation (22.1 vs.\ 12.5 tok/s
generation with 10 adapters). The base model alone runs at 60.8 tok/s;
sparse routing with adapters reaches 26.2 tok/s.

% PILOT: 10-adapter grove: 95/101 composition checks, generic PPL -14.2%
% PILOT: Sparse routing: +77% throughput vs dense (22.1 vs 12.5 tok/s)

% ----------------------------------------------------------------------------
\subsection{Knowledge PPL as Evaluation Metric}
\label{sec:knowledge-ppl}

Uniform perplexity underestimates knowledge injection. An adapter that
dramatically improves prediction of rare domain terms while slightly
degrading common tokens will show a modest uniform PPL change, because
common tokens dominate the average.

We split evaluation tokens by base-model surprisal into two populations:
\begin{itemize}[nosep,leftmargin=1.5em]
\item \textbf{Knowledge tokens}: tokens where $\ell_{\text{base}}(t)$
      exceeds the median---the tokens the base model finds surprising.
      These are where domain knowledge lives.
\item \textbf{Style tokens}: tokens where $\ell_{\text{base}}(t)$ is
      below the median---predictable tokens that reflect formatting,
      syntax, and common patterns.
\end{itemize}

% PILOT: Knowledge PPL -64.9% on knowledge tokens, +11.5% on style tokens

Pilot results show the adapter improves knowledge-token perplexity by
64.9\% while increasing style-token perplexity by 11.5\%. Uniform PPL
averages these into a modest improvement, hiding the large knowledge
gain.

% ============================================================================
\section{Experiments}
\label{sec:experiments}

All experiments use Qwen3-8B (36 layers, 8.2B parameters) as the base
model. Adapters use LoRA on FFN layers 12--35 (rank-16 unless stated
otherwise). We report means $\pm$ 95\% confidence intervals across seeds.

% PLACEHOLDER: Experiment sections below describe what will be measured.
% Tables contain [TBD] cells to be filled by publication experiment scripts.

% ----------------------------------------------------------------------------
\subsection{E1: Gate Protection}
\label{sec:e1}

\textbf{Hypothesis.} Per-layer delta gates protect generic capability
during domain adaptation, unlike standard PEFT LoRA which uniformly
applies adapter modifications.

\textbf{Setup.} Three conditions: (1) base model (no adapter),
(2) standard PEFT LoRA (rank-16, layers 12--35, no gate),
(3) gated adapter (same rank and layers, per-layer delta gate with
two-phase training). Training data: PubMed 50K documents. Evaluation:
MMLU, ARC-Challenge, HellaSwag, MedQA; plus domain PPL and generic PPL
on held-out C4. Seeds: 5.

% PILOT: PEFT +25-32% generic degradation, gated -6-13%
% PILOT: 4 seeds, consistent across LoRA implementations (custom Expert and standard PEFT)
% PILOT: Even naive output gate reduces degradation 11x (from +25% to +2.9%)

\begin{table}[h]
\centering
\caption{E1: Generic capability preservation during domain adaptation.
Mean $\pm$ 95\% CI across 5 seeds.}
\label{tab:e1}
\begin{tabular}{@{}lccccc@{}}
\toprule
Condition & Generic PPL $\Delta$ & MMLU & ARC-C & HellaSwag & MedQA \\
\midrule
Base (no adapter) & --- & [TBD] & [TBD] & [TBD] & [TBD] \\
PEFT LoRA          & [TBD] & [TBD] & [TBD] & [TBD] & [TBD] \\
Gated adapter      & [TBD] & [TBD] & [TBD] & [TBD] & [TBD] \\
\bottomrule
\end{tabular}
\end{table}

% ----------------------------------------------------------------------------
\subsection{E2: Multi-Adapter Composition}
\label{sec:e2}

\textbf{Hypothesis.} Independently-trained gated adapters compose
without catastrophic interference when gates are jointly fine-tuned.

\textbf{Setup.} Train 5 adapters independently on PubMed, The Stack
(Python), ArXiv (physics), ArXiv (CS), and BBC News. Deploy all 5
simultaneously with joint gate fine-tuning. Evaluate per-domain PPL,
cross-domain interference (PPL on domain $i$ with only adapter $j$
active), and generic PPL. Seeds: 3.

% PILOT: 10-adapter 95/101 checks, ARC +1%, HellaSwag +2.1%
% PILOT: Generic PPL -14.2% with 10 adapters
% PILOT: 6 failures are semantically justified overlaps

\begin{table}[h]
\centering
\caption{E2: Multi-adapter composition. Per-domain PPL improvement and
cross-domain interference. Mean $\pm$ 95\% CI across 3 seeds.}
\label{tab:e2}
\begin{tabular}{@{}lcccc@{}}
\toprule
Domain & Own adapter $\Delta$PPL & Best other $\Delta$PPL & Interference & Generic $\Delta$PPL \\
\midrule
PubMed     & [TBD] & [TBD] & [TBD] & [TBD] \\
Stack (Py) & [TBD] & [TBD] & [TBD] & [TBD] \\
ArXiv (ph) & [TBD] & [TBD] & [TBD] & [TBD] \\
ArXiv (CS) & [TBD] & [TBD] & [TBD] & [TBD] \\
BBC News   & [TBD] & [TBD] & [TBD] & [TBD] \\
\midrule
Combined   & [TBD] & --- & [TBD] & [TBD] \\
\bottomrule
\end{tabular}
\end{table}

% ----------------------------------------------------------------------------
\subsection{E3: Data Quality Ablation}
\label{sec:e3}

\textbf{Hypothesis.} Data quality, not architecture, is the bottleneck
for knowledge injection. The gate protects generic capability regardless
of data quality, but domain improvement requires curated data.

\textbf{Setup.} Train gated adapters on four medical data sources of
varying quality: (1) PubMed abstracts (curated), (2) FineWeb-Edu
keyword-filtered medical texts, (3) Wikipedia medical articles,
(4) Reddit medical discussions. Evaluate on MedQA and generic PPL.
Seeds: 3.

% PILOT: FineWeb keyword-filtered -> MedQA -4pp (degrades)
% PILOT: Curated BBC -> domain -17%, but keyword BBC -> much worse
% PILOT: Gate protects generic PPL regardless of data quality

\begin{table}[h]
\centering
\caption{E3: Data quality ablation for medical domain. MedQA accuracy
and generic PPL change. Mean $\pm$ 95\% CI across 3 seeds.}
\label{tab:e3}
\begin{tabular}{@{}lcc@{}}
\toprule
Data source & MedQA $\Delta$ & Generic PPL $\Delta$ \\
\midrule
PubMed abstracts      & [TBD] & [TBD] \\
FineWeb-Edu filtered  & [TBD] & [TBD] \\
Wikipedia medical     & [TBD] & [TBD] \\
Reddit medical        & [TBD] & [TBD] \\
\bottomrule
\end{tabular}
\end{table}

% ----------------------------------------------------------------------------
\subsection{E4: Gate Component Ablation}
\label{sec:e4}

\textbf{Hypothesis.} The per-layer delta gate is the critical component;
alternative gating strategies (output-level gate, uniform gate, no gate)
are inferior on at least one of generic PPL or domain PPL.

\textbf{Setup.} Six conditions: (1) no gate (standard LoRA),
(2) output-level scalar gate, (3) per-layer uniform gate (same value
all layers), (4) per-layer delta gate (ours), (5) per-layer gate
without delta (gates adapter output, not delta), (6) per-layer delta
gate with L1 sparsity. All use rank-16 on PubMed. Seeds: 3.

% PILOT: 8 combinations tested, per-layer delta gate best on both dimensions
% PILOT: L1 not special: L1 ~ L2, any sparsity pressure works

\begin{table}[h]
\centering
\caption{E4: Gate component ablation. Domain and generic PPL change
relative to base. Mean $\pm$ 95\% CI across 3 seeds.}
\label{tab:e4}
\begin{tabular}{@{}lcc@{}}
\toprule
Gate configuration & Domain PPL $\Delta$ & Generic PPL $\Delta$ \\
\midrule
No gate (standard LoRA)     & [TBD] & [TBD] \\
Output-level scalar         & [TBD] & [TBD] \\
Per-layer uniform           & [TBD] & [TBD] \\
Per-layer delta (ours)      & [TBD] & [TBD] \\
Per-layer non-delta         & [TBD] & [TBD] \\
Per-layer delta + L1        & [TBD] & [TBD] \\
\bottomrule
\end{tabular}
\end{table}

% ----------------------------------------------------------------------------
\subsection{E5: Active Learning Comparison}
\label{sec:e5}

\textbf{Hypothesis.} Active data selection by compression progress
outperforms random selection, and the combination of active selection
with gated adapters is strictly better than either alone.

\textbf{Setup.} Three conditions: (1) random batch selection + no gate,
(2) active selection (by compression progress) + no gate,
(3) active selection + gated adapter. Data pool: PubMed + Stack + ArXiv.
Evaluation: primary domain PPL and generic PPL after matched compute
budget. Seeds: 3.

% PILOT: Active+gates generic -3.9% vs random +7.5%
% PILOT: Active selection beats random by 6.4pp on primary domain

\begin{table}[h]
\centering
\caption{E5: Active learning comparison. PPL change after matched
compute budget. Mean $\pm$ 95\% CI across 3 seeds.}
\label{tab:e5}
\begin{tabular}{@{}lcc@{}}
\toprule
Condition & Domain PPL $\Delta$ & Generic PPL $\Delta$ \\
\midrule
Random + no gate       & [TBD] & [TBD] \\
Active + no gate       & [TBD] & [TBD] \\
Active + gated (ours)  & [TBD] & [TBD] \\
\bottomrule
\end{tabular}
\end{table}

% ----------------------------------------------------------------------------
\subsection{E6: Adapter Rank Pruning}
\label{sec:e6}

\textbf{Hypothesis.} Post-training SVD truncation of adapter weights
reveals per-layer rank requirements. Layers where rank-4 matches rank-16
quality can be pruned to 4$\times$ less compute, analogous to biological
synaptic pruning: assimilation layers (Piaget's lens) need low rank,
accommodation layers (crystal) need high rank.

\textbf{Setup.} Train a rank-16 gated adapter on PubMed. For each of
layers 12--35, truncate via SVD to rank-8, rank-4, rank-2, rank-1.
Measure per-layer PPL impact of truncation. Report the rank-quality
curve and identify which layers are rank-compressible. Seeds: 3.

% PILOT: Rank-selectivity curve peaks at rank-64 for 8B experiments
% PILOT: 5 data points at rank-64 on Qwen3-8B L12

\begin{table}[h]
\centering
\caption{E6: Rank pruning. Per-layer PPL impact of SVD truncation from
rank-16. Showing 6 representative layers. Mean $\pm$ 95\% CI across 3 seeds.}
\label{tab:e6}
\begin{tabular}{@{}lccccc@{}}
\toprule
Layer & Rank-8 $\Delta$ & Rank-4 $\Delta$ & Rank-2 $\Delta$ & Rank-1 $\Delta$ & Gate mag. \\
\midrule
L14 & [TBD] & [TBD] & [TBD] & [TBD] & [TBD] \\
L18 & [TBD] & [TBD] & [TBD] & [TBD] & [TBD] \\
L22 & [TBD] & [TBD] & [TBD] & [TBD] & [TBD] \\
L26 & [TBD] & [TBD] & [TBD] & [TBD] & [TBD] \\
L30 & [TBD] & [TBD] & [TBD] & [TBD] & [TBD] \\
L34 & [TBD] & [TBD] & [TBD] & [TBD] & [TBD] \\
\bottomrule
\end{tabular}
\end{table}

% ============================================================================
\section{Related Work}
\label{sec:related}

\paragraph{Selective training on tokens.}
Rho-1~\cite{rho1_2024} demonstrated that training on only 3\% of
pretraining tokens---selected by excess loss over a reference
model---improves math reasoning by 30\% absolute. DSIR~\citep{xie2023dsir}
achieves expert-curation-level performance via importance resampling in
n-gram feature space. Both are offline, static selection methods. Our
architecture enables \emph{online} selection: the adapter's own
compression progress determines what to learn next, continuously
adapting as the adapter's knowledge changes. The frozen base provides
the reference model that Rho-1 requires, without needing a separate
model.

\paragraph{Compression progress and curiosity.}
Schmidhuber~\cite{schmidhuber2008} formalized curiosity as the first
derivative of compression ability, arguing that agents should seek data
that maximizes their rate of learning---not data that maximizes surprise.
Pathak et al.~\cite{pathak2017icm} operationalized this in RL with the
Intrinsic Curiosity Module, filtering prediction error through a learned
feature space to separate learnable variation from noise. Our relative
learntropy (Equation~\ref{eq:relative-learntropy}) implements a similar
filter: the base model's loss provides the noise floor, and the excess
is the learnable signal.

\paragraph{LoRA and parameter-efficient fine-tuning.}
Hu et al.~\cite{hu2022lora} showed that fine-tuning updates lie in a
low-rank subspace. Biderman et al.~\cite{biderman2024lora} found that
LoRA learns less than full fine-tuning but also forgets less---a
trade-off that our gating mechanism exploits. We extend standard LoRA
with per-layer delta gating: rather than applying all adapter layers
uniformly, the gate learns selective activation from data. HydraLoRA~\cite{hydralora2024} uses asymmetric adapters for multi-task
learning; MoLoRA~\cite{molora2026} adds per-token routing. Our delta
gate operates at a different level: per-layer, per-adapter, controlling
\emph{whether} to apply the adapter at each layer rather than
\emph{which} adapter to apply.

\paragraph{Mixture-of-Experts.}
Standard MoE architectures~\cite{mixtral2024,deepseek2024moe,deepseekv3}
route tokens to experts via learned gates with load-balancing losses.
This produces functional specialization (structure vs.\ content) but
resists domain specialization~\cite{standing2026,openmoe2024,domaindriver2026}.
DeepSeek-V3~\cite{deepseekv3} introduces shared always-on experts
alongside routed specialists---a design our trunk-adapter split
mirrors, with the frozen base as the permanent shared expert.
DynMoE~\cite{l2r2026} allows dynamic expert count; our speculative
splitting with rollback achieves similar flexibility while being
data-driven rather than architecture-driven.

\paragraph{Layer efficiency.}
Recent work shows that many transformer layers are redundant: ShortGPT
removes up to 25\% of layers with minimal quality loss; Mixture-of-Depths
(MoD) routes tokens to skip layers dynamically. Our per-layer gate
discovers a similar structure from data: layers with near-zero gate
activation contribute nothing for a given domain and could be skipped
entirely.

% ============================================================================
\section{Limitations}
\label{sec:limitations}

\paragraph{Single model family.}
All experiments use Qwen3-8B. The delta gate mechanism is
architecture-agnostic (it requires only FFN layers with residual
connections), but we have not validated on other model families.
Generalization to Llama, Mistral, or encoder-decoder architectures
remains untested.

\paragraph{Data quality is the bottleneck.}
The gate protects generic capability regardless of training data quality,
but domain improvement requires high-quality, curated data. Pilot
experiments show that keyword-filtered web data (FineWeb-Edu)
\emph{degrades} medical benchmarks by 4--10 percentage points despite
the gate's protection. Architecture cannot compensate for data.

\paragraph{Hallucination risk.}
The delta gate detects \emph{domain membership}---whether the input
matches the adapter's training distribution---not \emph{factual
correctness}. An adapter confidently activated on a medical query may
still hallucinate. The gate is a necessary but not sufficient
condition for reliable domain knowledge.

\paragraph{Active learning not yet tested at scale.}
The compression progress signal and active learning cycle
(Section~\ref{sec:active-cycle}) have been validated in components (token
selection, Goldilocks zone detection) but not as an integrated pipeline
on large data pools. Scaling from pilot experiments (500 documents) to
publication-scale data (50K+ documents) may reveal new failure modes.

\paragraph{No negative results from scaled experiments yet.}
The pilot results in this paper come from experiments on 500--5000
documents. We report them honestly as pilots, not as definitive findings.
The publication experiments (Section~\ref{sec:experiments}) will provide
the necessary statistical power.

% ============================================================================
\section{Conclusion}
\label{sec:conclusion}

We have presented the \grove{}, a framework for placing the long tail
of language model knowledge in self-organizing, hot-pluggable adapters.
Three insights drive the design:

\begin{enumerate}[nosep,leftmargin=1.5em]
\item \textbf{The frozen base as epistemic reference.} By never modifying
      base weights, we gain a permanent reference point for measuring
      what any adapter has learned. This is not a constraint---it is the
      key enabler of relative learntropy, compression progress tracking,
      and the active learning cycle.
\item \textbf{Per-layer delta gates as regime selectors.} The gate
      implements a Piagetian spectrum from assimilation (base sufficient,
      gate near zero) to accommodation (adapter overrides base, gate
      near one). This protects generic capability while enabling
      selective knowledge injection---a trade-off that standard PEFT
      LoRA cannot achieve.
\item \textbf{Knowledge PPL reveals what uniform PPL hides.} Splitting
      evaluation by base-model surprisal shows that adapters improve
      knowledge-token perplexity by 64.9\%, a gain that uniform PPL
      averages away.
\end{enumerate}

The grove is designed for community-scale growth: each adapter trains
independently on commodity hardware, requiring only the frozen base and
a domain-specific data stream. Contributors add specialists without
retraining the core. The long tail of human knowledge is long---no
single organization can curate it all. A self-organizing, composable
system of lightweight specialists may be the natural structure for
capturing it.

\paragraph{Future directions.}
Three extensions follow naturally: (1)~\emph{layer skipping}---layers
with near-zero gate activation can be skipped entirely at inference time,
connecting to Mixture-of-Depths research; (2)~\emph{curious
adapters}---using compression progress not just for token selection but
for autonomous exploration of new data domains; and (3)~\emph{scaled
active learning}---validating the full active learning cycle on large,
diverse data pools where the grove must discover its own curriculum
without human curation.

% ============================================================================
\begin{thebibliography}{99}

\bibitem{hu2022lora} E. Hu et al. LoRA: Low-Rank Adaptation of Large Language Models. \textit{ICLR}, 2022.

\bibitem{geva2021transformer} M. Geva, R. Schuster, J. Berant, O. Levy. Transformer Feed-Forward Layers Are Key-Value Memories. \textit{EMNLP}, 2021.

\bibitem{li2018intrinsic} C. Li et al. Measuring the Intrinsic Dimension of Objective Landscapes. \textit{ICLR}, 2018.

\bibitem{aghajanyan2021intrinsic} A. Aghajanyan et al. Intrinsic Dimensionality Explains the Effectiveness of Language Model Fine-Tuning. \textit{ACL}, 2021.

\bibitem{komatsuzaki2023sparse} A. Komatsuzaki et al. Sparse Upcycling. \textit{arXiv:2212.05055}, 2023.

\bibitem{bena2025modularity} G. B\'{e}na, D. Goodman. Dynamics of Specialization Under Resource Constraints. \textit{Nature Communications}, 2025.

\bibitem{l2r2026} L2R: Low-Rank Lipschitz-Controlled Routing. \textit{arXiv:2601.21349}, 2026.

\bibitem{shannon1948} C. E. Shannon. A Mathematical Theory of Communication. \textit{Bell System Technical Journal}, 27(3):379--423, 1948.

\bibitem{hartigan1985dip} J. A. Hartigan, P. M. Hartigan. The Dip Test of Unimodality. \textit{Annals of Statistics}, 13(1):70--84, 1985.

\bibitem{rho1_2024} Z. Lin et al. Rho-1: Not All Tokens Are What You Need. \textit{arXiv:2404.07965}, 2024.

\bibitem{schmidhuber2008} J. Schmidhuber. Driven by Compression Progress: A Simple Principle Explains Essential Aspects of Subjective Beauty, Novelty, Surprise, Interestingness, Attention, Curiosity, Creativity, Art, Science, Music, Jokes. \textit{arXiv:0812.4360}, 2008.

\bibitem{piaget1952} J. Piaget. \textit{The Origins of Intelligence in Children}. International Universities Press, 1952.

\bibitem{vygotsky1978} L. S. Vygotsky. \textit{Mind in Society: The Development of Higher Psychological Processes}. Harvard University Press, 1978.

\bibitem{wozniak2018learntropy} P. A. Wozniak. Learntropy. \textit{SuperMemo Guru}, 2018. \url{https://supermemo.guru/wiki/Learntropy}.

\bibitem{standing2026} The Illusion of Specialization: Standing Committee in Mixture-of-Experts. \textit{arXiv:2601.03425}, 2026.

\bibitem{deepseek2024moe} X. Dai et al. DeepSeekMoE: Towards Ultimate Expert Specialization. \textit{ACL}, 2024.

\bibitem{deepseekv3} DeepSeek-AI. DeepSeek-V3 Technical Report. \textit{arXiv:2412.19437}, 2024.

\bibitem{hydralora2024} T. Tian et al. HydraLoRA: An Asymmetric LoRA Architecture for Efficient Fine-Tuning. \textit{NeurIPS}, 2024.

\bibitem{molora2026} MoLoRA: Composable Specialization via Per-Token Adapter Routing. \textit{arXiv:2603.15965}, 2026.

\bibitem{openmoe2024} F. Xue et al. OpenMoE: An Early Effort on Open Mixture-of-Experts Language Models. \textit{ICML}, 2024.

\bibitem{mixtral2024} A. Jiang et al. Mixtral of Experts. \textit{arXiv:2401.04088}, 2024.

\bibitem{domaindriver2026} What Gets Activated: Uncovering Domain and Driver Experts in MoE Language Models. \textit{arXiv:2601.10159}, 2026.

\bibitem{jang2017gumbel} E. Jang, S. Gu, B. Poole. Categorical Reparameterization with Gumbel-Softmax. \textit{ICLR}, 2017.

\bibitem{biderman2024lora} S. Biderman et al. LoRA Learns Less and Forgets Less. \textit{TMLR}, 2024.

\bibitem{pathak2017icm} D. Pathak et al. Curiosity-driven Exploration by Self-Supervised Prediction. \textit{ICML}, 2017.

\bibitem{xie2023dsir} S. M. Xie et al. Data Selection with Importance Resampling. \textit{NeurIPS}, 2023.

\end{thebibliography}

\end{document}
